\chapter{Kinematics of a Continuum}

\section{Terminologies and definitions}

\begin{itemize}
\item \fbox{Body}\\

A body occupies regions in $E_3$ and is composed of material points(or particles) to each of which the density of mass, momentum, and energy etc. are defined throughout the body in a continuously except, possibly, for a finite number of singular points, lines, surface, or volumes of discontinuously such as holes, cracks, voids, and cavities in the body which require special attention. 

\item \fbox{Kinematics}\\

Study of motion of bodies regardless the nature of force that cause the motion.

\item \fbox{Deformations}\\

Let a body occupies a region $R\in E_3$. Every particle can be located by its position vector $\tenq{x}$ from some fixed point $O$. When the body is deformed, every particle takes up new position vector denoted by $\tenq{y}$. (referred to the same origin $O$)

\begin{figure}[H]
\centering
\resizebox{0.6\textwidth}{!}{\input{figure3.eps_tex}}

\caption{Scheme for $R$}
\label{Scheme for $R$ 3}
\end{figure}

Then the deformation of the body is known if the function 

\begin{equation}
\tenq{y} = \hat{\tenq{y}}(\tenq{x}) \text{ or }y_i = \hat{y}_i(\tenq{x}) \label{eq:3-1}\tag{1}
\end{equation}

is known. 

\end{itemize}

\begin{note}\mbox{}\\

\begin{enumerate}[label=(\arabic*)]

\item
It is possible that the body can occupy the same in two different states in which particles occupy different positions. 
\begin{enumerate}[label= e.g.\alph*)]
\item 
A body occupies a sphere in $E_3$. It is then rotated but occupies the same sphere. The two state should be distinguishable. (an example of \emph{rigid deformation})

\item 
$R$ is a circular cylinder. Twist a rubble rod occupying $R$ and it is will still do so.
\end{enumerate}

\item 
In continuum mechanics, we demand the mapping \eqref{eq:3-1} be 

\begin{enumerate}[label= (\alph*)]
\item 
1 to 1 (globally), i.e. $\tenq{y} = \hat{\tenq{y}}(\tenq{x})$, $\tenq{x} = \hat{\tenq{x}}(\tenq{y})$ and $\hat{\tenq{y}}^{-1}(\tenq{x}) = \hat{\tenq{x}}(\tenq{y}) = \tenq{x} \phantom{-} \forall \tenq{x}\in R$ \label{deformation-a}

\item 
$\hat{\tenq{y}} \in C^1(R)$ and $\hat{\tenq{y}}^{-1} \in C^1(R')$\label{deformation-b}
\item 
$\det\nabla \hat{\tenq{y}}>0$ on $R$\label{deformation-c}

\end{enumerate}

Property \ref{deformation-a} guarantees that the body not penetrate itself, i.e. distinct particles occupy distinct positions in every state. 

Property \ref{deformation-b} assumes \emph{no tearing}.

Property \ref{deformation-c} represents two things: 
\hfill
\begin{itemize}

\item
1 to 1 correspondence of local directions.(curves that intersect each other at nonzero angle are not mapped on to the tangent ones)

\item
locally the volume of a part of the body(a subregion) per unit original volume is always grater than zero. 

\end{itemize}

Note that properties \ref{deformation-b}+\ref{deformation-c} does \emph{not} imply that $\hat{\tenq{y}}$ is globally 1 to 1 on $R$. They only imply that $\hat{\tenq{y}}$ is locally 1 to 1, i.e. could have $\det\nabla \hat{\tenq{y}}>0$ with $\tenq{y} = \hat{\tenq{y}}(\tenq{x}) = \hat{\tenq{y}}(\tenq{x'})$. 

\begin{figure}[H]
\centering
\resizebox{0.6\textwidth}{!}{\input{figure4.eps_tex}}

\caption{Scheme for properties \ref{deformation-b}+\ref{deformation-c}}
\label{Scheme for properties}
\end{figure}

\item 

$\tenq{y} = \hat{\tenq{y}}(\tenq{x})$, $\tenq{x} = \hat{\tenq{x}}(\tenq{y})$ and $\hat{\tenq{x}} = \hat{\tenq{y}}^{-1}\colon R' \to R$\\
Then 
\[\hat{\tenq{y}} ( \hat{\tenq{x}}(\tenq{y}) )= \tenq{y} \phantom{-} \forall \tenq{y} \in R'\]
By chain rule
\[
\frac{\partial \hat{y}_i(\tenq{x})}{\partial x_k}\frac{\partial \hat{x}_k}{\partial y_j} = \frac{\partial y_i}{\partial y_j} = \delta_{ij}
\]

\begin{equation}
\Rightarrow \boxed{\nabla \hat{\tenq{y}}(\tenq{x})\nabla \hat{\tenq{x}}(\tenq{y}) = \tenq{1}}\Rightarrow \det\nabla \hat{\tenq{y}}\det\nabla \hat{\tenq{x}} = 1  \label{eq:det of deformation}\tag{$\bigstar$}
\end{equation}

or

\[\det\nabla \hat{\tenq{x}}(\tenq{y}) = \frac{1}{\det\nabla \hat{\tenq{y}}(\tenq{x})}\phantom{-} \forall \tenq{x}\in R\]

\end{enumerate}
\end{note}

\begin{definition}
We say that $\hat{\tenq{y}}\colon R \to R'$ is an \emph{admissible deformation} if 

\begin{enumerate}[label= (\alph*)]

\item
$\hat{\tenq{y}}$ is 1 to 1 on $R$
\item 
$\hat{\tenq{y}} \in C^2(R)$
\item 
$\det\nabla \hat{\tenq{y}} > 0$ on $R$
\end{enumerate}
\begin{note}
By the \emph{smooth inverse theorem}\footnotemark, then $\hat{\tenq{x}} = \hat{\tenq{y}}^{-1}$ on $R'$ is also admissible, i.e. $\hat{\tenq{x}} \in C^2(R')$; while $\det\nabla \hat{\tenq{x}}>0$ on $R'$ by equation \eqref{eq:det of deformation}. 
\end{note}
\footnotetext{referred to advanced calculus}
\end{definition}

\section{Motion of a body}

For fixed choice of the region $R$ (called \emph{reference configuration})

\[
\hat{\tenq{y}}(\tenq{x},x) \colon R \to R_t, t\in [t_0,t_1]\equiv\tau
\]
represent the motion of the body.

\begin{figure}[H]
\centering
\resizebox{0.8\textwidth}{!}{\input{figure5.eps_tex}}

\caption{the motion of a body}
\label{Motion of a body}
\end{figure}

\begin{note}
The reference configuration $R$ is arbitrary. It needs not occupied by the body at any time during the motion. You may even choose \emph{fictitious} one if it can help to simplify the formulation of the problem. The technique is often used in plastic mechanics. 
\end{note}
A continuous motion is one such that 
\[\tenq{y} = \hat{\tenq{y}}(\bullet,t) \in C(\tau)\]
\begin{definition}
The Jacobian of the motion $\hat{\tenq{y}}(\bullet,\bullet)$ on $R\times \tau$ is 
\[J(\tenq{x},t) = \det\nabla\hat{\tenq{y}}(\tenq{x},t)\phantom{-}\forall (\tenq{x},t)\in R \times \tau\]
\end{definition}

If we choose $R_{t*}$ as reference configuration where $t_* \in [t_0,t_1]$ then 

\[\hat{\tenq{y}}(\tenq{x},t_*) = \tenq{x}\]

\[\Rightarrow \nabla \hat{\tenq{y}}(\tenq{x},t_*) = \tenq{1}\Rightarrow J(\tenq{x},t_*) = 1\]
\[J(\tenq{x},t) > 0 \phantom{-}\forall (\tenq{x},t)\in R \times \tau\]

\begin{figure}[H]
\centering
\resizebox{0.6\textwidth}{!}{\input{figure6.eps_tex}}

\caption{Jacobian versus time}
\label{jacobian}
\end{figure}

\begin{definition}

A motion is admissible if $\hat{\tenq{y}} \in C^2(R \times \tau)$

\end{definition}


\begin{note}\mbox{}\\

\begin{enumerate}[label=(\alph*)]

\item 
Admissible motions are smooth one-parameter families of admissible deformations, $\hat{\tenq{y}}(\tenq{x},\bullet)$.

\item
$\frac{\partial \hat{\tenq{y}}(\tenq{x},t)}{\partial t},\frac{\partial^2 \hat{\tenq{y}}(\tenq{x},t)}{\partial t^2}$ exist and are continuous $\forall (\tenq{x},t)\in R \times \tau$.

\item 
There exists $\tenq{y}^{-1}(\bullet,t) = \tenq{x}(\bullet,t)$ defined on $R_t \phantom{-} \forall t \in \tau$. 

\[\hat{\tenq{x}}(\bullet,t)\colon R_t \to R\]
\[\hat{\tenq{x}}(\hat{\tenq{y}}(\tenq{x},t),t) = \tenq{x}\]
\[\hat{\tenq{y}}(\hat{\tenq{x}}(\tenq{y},t),t) = \tenq{y}\]
\[\hat{\tenq{x}}(\bullet,t)\in C^2(R_t)\phantom{-} \forall t \in \tau\]
$\Rightarrow \frac{\partial \hat{\tenq{x}}}{\partial t},\frac{\partial^2 \hat{\tenq{x}}}{\partial t^2}$ exist and are continuous on $m \buildrel \Delta \over = \left\{ (\tenq{y},t),y\in R_t,t\in \tau\right\}$.

\item
Let $P$ be the particle of the body and $\tenq{y} = \hat{\tenq{y}}(\tenq{x},t)$ is the position vector of $P$ at time. 

\begin{alignat*}{3}
& x_i   \quad && \cdots  \quad && \text{Lagrangian coordinate of }P\\
& y_i   \quad && \cdots  \quad && \text{Eulerian coordinate of }P
\end{alignat*}

For fixed $\tenq{x}$, $\hat{\tenq{y}}(\tenq{x},t)$ is the path line of $P$.

\item 
Given a motion let $\Psi$ be a field defined on $R\times\tau$, $\Psi = \Psi(\tenq{x},t), \tenq{x}\in R, t\in \tau$, then one may define $\psi$ on $m$ by 

\[\psi(\tenq{y},t) = \Psi(\hat{\tenq{x}}(\tenq{y},t),t)\]
where
\[\hat{\tenq{x}}(\bullet,t) = \tenq{y}^{-1}(\bullet,t)\]

Equivalently, given $\psi(\tenq{y},t)$ on $m$ one can define $\Psi (\tenq{x},t)$ such that 
\[\Psi(\tenq{x},t) = \psi(\hat{\tenq{y}}(\tenq{x},t),t)\phantom{-} \forall (\tenq{x},t) \in R \times \tau\]
In this case we refer to $\Psi (\tenq{x},t)$ as a reference (or a material) field and $\psi (\tenq{y},t)$ as spatial field. 
\end{enumerate}
\end{note}

\begin{example}
Consider the motion of a body

\begin{align*}
y_1 &= x_1 + ktx_2\\
y_2 &= x_2\\
y_3 &= x_3
\end{align*}

The Lagrangian coordinate $\tenq{x}$ gives the position of a particle at $t = 0$.

\begin{enumerate}[label=(\alph*)]

\item 
Sketch the configuration of a unit cube at time $t$. 
\item 
If the temperature field is given by 

\[\theta = y_1 + y_2\]
find the material field of temperature.

\item 
Find the velocity the rate of change of temperature for particular material particles. Express both in material and spatial description.

\end{enumerate}
\end{example}

\begin{solution}\mbox{}\\

\begin{figure}[H]
\centering
\resizebox{0.4\textwidth}{!}{\input{figure7.eps_tex}}

\caption{Scheme of example}
\label{example L/E}
\end{figure}

\begin{enumerate}[label=(\alph*)]

\item 
By particle on $OA$
\begin{align*}
\left(x_1,x_2,x_3\right)_{OA} &= \left(x_1,0,x_3\right)\\
\left(y_1,y_2,y_3\right)_{OA} &= \left(x_1,0,x_3\right)
\end{align*}
i.e. $OA$ does not change.\\

For $CB$
\begin{align*}
\left(x_1,x_2,x_3\right)_{CB} &= \left(x_1,1,x_3\right)\\
\left(y_1,y_2,y_3\right)_{CB} &= \left(x_1+kt,1,x_3\right)
\end{align*}
i.e. $CB$ move horizontally through a distance $kt$.\\

For $OC$
\begin{align*}
\left(x_1,x_2,x_3\right)_{OC} &= \left(0,x_2,x_3\right)\\
\left(y_1,y_2,y_3\right)_{OC} &= \left(ktx_2,x_2,x_3\right)
\end{align*}
i.e. particle on $OC$ move horizontally through a distance linearly proportional to its height.\\
Similarly for $AB$.

\item
$\theta = y_1 + y_2 = x_1 + ktx_2 + x_2 = x_1 + \left(kt+1\right)x_2$

\item Velocity components of the particle with Lagrangian position vector $\tenq{x}$. 
\[{v_i} = {\left. {\frac{{\partial {y_i}(\tenq{x},t)}}{{\partial t}}} \right|_{\tenq{x} \text{ fixed}}}\]

\[\Rightarrow\left\{\begin{aligned}
& v_1 = kx_2,v_2 = 0,v_3 = 0 \cdots \text{material field of velocity}\\
& v_1 = ky_2,v_2 = 0,v_3 = 0 \cdots \text{spatial field of velocity}\end{aligned}\right.\]

Rate of change of temperature

\begin{align*}
{\left. {\frac{{\partial \theta }}{{\partial t}}} \right|_{\tenq{x} \text{ fixed}}} &= kx_2 \cdots \text{material field of temperature}\\
&= ky_2 \cdots \text{spatial field of temperature}
\end{align*}

\end{enumerate}

\end{solution}